\subsection*{Пояснение к разложению.}
Итак, мы получили формулу вида
\[
\mathbb{E}[(y - \hat{y})^2] = \text{Bias}^2 + \text{Variance} + \text{Irreducible Error}
\]

Давайте разбирать компоненты по очереди.
\subsubsection*{Смещение}
\[
Bias = \mathbb{E}[\hat{y}] - y
\]

То есть, насколько в среднем наши прогнозы отличаются от реальных ответов.
Ещё смещение иногда называют \textit{систематической ошибкой}. Величина означает ошибку, которую модель допускает систематически: неважно, какой размер данных или какого качества данные "--- модель в любом случае будет ошибаться.

Обычно смещение связывают с простотой модели: если модель слишком простая (представьте себе дерево глубины 1 или 2), то она гарантированно будет \textit{смещена} (т.е. допускать эту систематическую ошибку).

Представьте, что вы строите модель предсказания цены квартиры. Истинная закономерность вот такая:
\[
y = 50 * S + 500 * R + 100 * A + \epsilon
\]
т.е. некоторая линейная зависимость между площадью, количеством комнат и возрастом здания.

А мы, не зная этого, построили слишком простую модель вида
\[
\hat{y} = 40 * S
\]

\newpage
Давайте посмотрим, что будет происходить на разных наборах данных:

\begin{enumerate}
    \item Квартира: $
    S = 60, R = 2, A = 10
    $
    
    Реальная цена: $
    y_1 = 50 * 60 + 500 * 2 + 100 * 10 = 5000
    $
    
    Предсказание нашей модели: $
    40 * 60 = 2400
    $
    
    Ошибка составляет $2600$

    \item Квартира: $
    S = 70, R = 3, A = 15
    $
    
    Реальная цена: $
    y_2 = 50 * 70 + 500 * 3 + 100 * 15 = 6500
    $
    
    Предсказание: $
    40 * 70 = 2800
    $
    
    Ошибка составляет $3700$

    \item Квартира: $
    S = 50, R = 1, A = 5
    $
    
    Реальная цена: $
    y_3 = 50*50 + 500*1 + 100*5 = 3500
    $
    
    Предсказание: $
    40 * 50 = 2000
    $
    
    Ошибка составляет $1500$
\end{enumerate}
Итак, среднее предсказание модели:
\[
\mathbb{E}[\hat{y}] = \frac{\hat{y_1} + \hat{y_2} + \hat{y_3}} {3} = 
\frac{7200}{3} = 2400
\]
А целевая переменная в среднем:
\[
\mathbb{E}[y] = \frac{y_1 + y_2 + y_3}{3} = \frac{15000}{3} = 5000
\]
Теперь по формуле считаем смещение:
\[
Bias = 2400 - 5000 = -2600
\]

Выходит, что построенная нами модель систематически недопредсказывает цену квартиры на $2600$ в среднем.

\subsubsection*{Разброс}
\[
Variance = \mathbb{E}[(\hat{y} - \mathbb{E}[\hat{y}])^2]
\]

Мы измеряем, насколько сильно одно конкретное предсказание отклоняется от всех усреднённых предсказаний. Интуитивно это означает, насколько модель чувствительна к изменениям в данных. Вспомните пример из \texttt{SVM}, когда мы говорили об уверенности классификатора. Если точка стоит рядом с разделяющей гиперплоскостью, то достаточно слегка подвинуть гиперплоскость, как класс на этом объекте изменится. Именно такая интуиция применима к понятию разброса.

В то время, как смещение связывают с простотой модели, разброс наоборот, со сложностью модели. Попытайтесь задуматься: если модель слишком сильно улавливает какие"=то зависимости в данных, то если расставить эти объекты на плоскости немного иным образом (т.е. поменять выборку на похожую), то результат сильно просядет. Это именно тот эффект, который мы наблюдаем при переобучении модели.

Давайте предположим, что по аналогии со смещением, у нас есть модель, которая трижды обучена на разных датасетах (квартира из центра города, из пригорода и смешанные). Их предсказания: $\hat{y_1}=5500, \hat{y_2}=4200, \hat{y_3}=4800$. Соответственно,
\[
\mathbb{E}[\hat{y}] = \frac{\hat{y_1} + \hat{y_2} + \hat{y_3}}{3} = \frac{14500}{3} \approx 4833
\]

Посчитаем отклонение каждого конкретного прогноза от среднего:
\[
\hat{y_1} - \mathbb{E}[\hat{y}] = 5500 - 4833 = 667
\]

\[
\hat{y_2} - \mathbb{E}[\hat{y}] = 4200 - 4833 = -633
\]

\[
\hat{y_3} = \mathbb{E}[\hat{y}] = 4800 - 4833 = -33
\]

Ну и квадраты отклонений:
\[
(\hat{y_1} - \mathbb{E}[\hat{y}])^2 = 444,889
\]

\[
\hat{y_2} - \mathbb{E}[\hat{y}] = 400,689
\]

\[
\hat{y_3} - \mathbb{E}[\hat{y}] = 1,089
\]

Таким образом, в среднем разброс составляет
\[
\frac{846,667}{3} \approx 282,222
\]

Как мы видим, для разных типов квартир (из центра, смешанные и т.д.) получаются сильно разные ответы. Это и есть разброс модели: небольшие изменения в наборе данных сильно влияют на работу модели, что похоже на переобучение модели.

Но какое нам дело, что там на других наборах данных? На нашем работает "--- ну и хорошо, так ведь? Нет, не так ведь. Дело в том, что в рамках задачи данные, например, могут поступать ежедневно: при том же прогнозе погоды. В таком случае крайне важно, чтобы модель не выучивала каждый набор данных, а искала какие"=то общие закономерности (это и есть обобщающая способность модели).

Внимательный читатель мог заметить интересный парадокс: выходит, если мы упрощаем модель, то ошибка увеличивается, но если мы усложняем модель "--- ошибка тоже увеличивается? На самом деле, так оно и есть, и вскоре мы в этом убедимся. В машинном обучении главное находить компромисс между ошибкой, которая возникает из"=за смещения, и ошибкой, которая возникает из"=за разброса. Этот компромисс ещё называют \texttt{Bias-variance tradeoff}.

\subsubsection*{Шум}

У шума нет как таковой формулы. Вспомните из ранее описанной формулировки, что из себя представлял $\epsilon$. Так вот, этот эпсилон и есть шум "--- это некоторая ошибка, на которую невозможно повлиять и её никак нельзя уменьшить. Резюмируя "--- любая модель будет допускать ошибку, которую невозможно исправить и модель будет гарантированно ошибаться. В наших силах лишь сделать эту ошибку как можно меньше, но убрать её полностью "--- невозможно.